\section{The central limit theorem and the shape of fluctuations}

The law of large numbers says that
\[
\overline X_n
\]
becomes close to
\[
\mu
\]
when the sample size is large. It explains stability, but it does not describe the
remaining error
\[
\overline X_n-\mu.
\]

To use an observed sample mean for inference, we need more information. We need
to know the typical size of this error and the approximate shape of its sampling
distribution.

The central limit theorem provides this second regularity.

\subsection*{The scale of the fluctuations}

Suppose
\[
X_1,X_2,\ldots,X_n
\]
are independent and identically distributed, with
\[
\E[X_i]=\mu
\qquad\text{and}\qquad
\Var(X_i)=\sigma^2<\infty.
\]
Then
\[
\E[\overline X_n]=\mu
\]
and
\[
\operatorname{SD}(\overline X_n)=\frac{\sigma}{\sqrt n}.
\]

Therefore the natural scale of the difference
\[
\overline X_n-\mu
\]
is not \(1\), and not \(1/n\), but
\[
\frac1{\sqrt n}.
\]

The standardized fluctuation is
\[
Z_n
=
\frac{\overline X_n-\mu}{\sigma/\sqrt n}.
\]
This quantity measures how many standard deviations the sample mean lies above
or below the population mean.

\subsection*{The central limit theorem}

Under the assumptions above, the distribution of \(Z_n\) approaches the standard
normal distribution as \(n\) increases:
\[
Z_n
=
\frac{\overline X_n-\mu}{\sigma/\sqrt n}
\xrightarrow{d}
N(0,1).
\]

\begin{theorembox}
For independent identically distributed observations with finite mean \(\mu\) and
finite positive variance \(\sigma^2\),
\[
\frac{\overline X_n-\mu}{\sigma/\sqrt n}
\xrightarrow{d}N(0,1).
\]
This statement is called the \textbf{central limit theorem}.
\end{theorembox}

The notation
\[
\xrightarrow{d}
\]
means convergence in distribution. For this course, the essential interpretation
is that the standardized sampling distribution becomes approximately normal for
large \(n\).

Equivalently, the unstandardized sample mean is approximately distributed as
\[
\overline X_n
\approx
N\left(\mu,\frac{\sigma^2}{n}\right).
\]
This is an approximation, not an exact identity in general.

\subsection*{What the theorem adds to the law of large numbers}

The two limit theorems answer different questions.

The law of large numbers says
\[
\overline X_n\approx\mu
\]
for large \(n\). It identifies the value around which the statistic stabilizes.

The central limit theorem says
\[
\overline X_n-\mu
\]
usually has size approximately
\[
\frac{\sigma}{\sqrt n},
\]
and that after dividing by this scale, the resulting distribution is
approximately standard normal.

Thus:
\[
\text{law of large numbers}
\quad\longrightarrow\quad
\text{location of the statistic},
\]
while
\[
\text{central limit theorem}
\quad\longrightarrow\quad
\text{shape and scale of its fluctuations}.
\]

\begin{summarybox}
The law of large numbers explains why the sample mean approaches \(\mu\). The
central limit theorem explains how the sample mean continues to fluctuate around
\(\mu\) when the sample size is large.
\end{summarybox}

\subsection*{The population need not be normal}

A common misunderstanding is that the central limit theorem requires the
individual observations to have a normal distribution. It does not.

The population distribution may be discrete, skewed, or otherwise non-normal.
The theorem concerns the sampling distribution of the standardized mean, not the
distribution of each individual observation.

For example, suppose
\[
X_i\sim\operatorname{Bernoulli}(p).
\]
Each \(X_i\) takes only the values \(0\) and \(1\), so its distribution is not
continuous and is certainly not normal. Nevertheless, the sample proportion
\[
\widehat p_n=\frac1n\sum_{i=1}^nX_i
\]
is a sample mean. Since
\[
\E[X_i]=p
\qquad\text{and}\qquad
\Var(X_i)=p(1-p),
\]
the central limit theorem gives
\[
\frac{\widehat p_n-p}
{\sqrt{p(1-p)/n}}
\approx N(0,1)
\]
for sufficiently large \(n\).

Equivalently,
\[
\widehat p_n
\approx
N\left(p,\frac{p(1-p)}{n}\right).
\]

\subsection*{A numerical interpretation}

Suppose a Bernoulli model has
\[
p=0.4
\]
and the sample size is
\[
n=400.
\]
Then
\[
\E[\widehat p]=0.4
\]
and
\[
\operatorname{SD}(\widehat p)
=
\sqrt{\frac{0.4(0.6)}{400}}
\approx 0.0245.
\]

The central limit theorem suggests that the standardized quantity
\[
\frac{\widehat p-0.4}{0.0245}
\]
is approximately standard normal.

A sample proportion such as
\[
\widehat p=0.42
\]
is less than one standard deviation above \(0.4\), so it is not surprising under
the model. A value such as
\[
\widehat p=0.50
\]
is more than four standard deviations above \(0.4\), so it would be highly
unusual under the same model.

This is the beginning of statistical inference. The observed statistic is judged
relative to the sampling distribution predicted by a proposed model.

\subsection*{When the approximation may be poor}

The phrase ``for large \(n\)'' does not refer to one universal sample size. The
quality of the approximation depends on the population distribution.

The normal approximation may be inaccurate when:

\begin{itemize}
    \item the sample is small;
    \item the population distribution is extremely skewed;
    \item very rare and very large observations strongly affect the mean;
    \item the variance is infinite;
    \item the observations are strongly dependent.
\end{itemize}

For Bernoulli data, the approximation is usually more reliable when both
\[
np
\qquad\text{and}\qquad
n(1-p)
\]
are not too small. These quantities represent the expected numbers of successes
and failures.

The central limit theorem is powerful because its conclusion is widely
applicable, but every application still depends on assumptions about the sampling
mechanism.

\begin{remarkbox}
A large sample does not repair every defective model. Limit theorems describe the
consequences of assumptions; they do not justify independence, identical
distributions, or good data collection.
\end{remarkbox}

\subsection*{From approximation to inference}

The approximation
\[
\frac{\overline X_n-\mu}{\sigma/\sqrt n}
\approx N(0,1)
\]
allows probabilities about the statistic to be computed using the standard
normal distribution.

It tells us which deviations from \(\mu\) are ordinary and which are unusual. It
also provides a way to construct ranges of parameter values that are compatible
with an observed statistic.

To use this idea systematically, we need a name for the standard deviation of a
sampling distribution. This quantity is called the standard error.

\begin{summarybox}
The central limit theorem provides a common approximate shape for many sampling
distributions. After centring by the parameter and scaling by the standard
deviation of the statistic, the remaining fluctuation is approximately standard
normal. This regularity turns sampling distributions into practical tools for
measuring uncertainty.
\end{summarybox}
